% Capítulo 5 — Transformaciones euclídeas y similaridades
\chapter{Transformaciones euclídeas y similaridades}\label{cap:05}\label{sec:bloque-euclideas}

\section{La idea}

Reflejamos y escalamos. Con reflexiones se construyen rotaciones, y con una escala común se pasa a las similaridades. Seguimos el hilo de las clases dictadas, con los resúmenes del profesor como guía.

\subsection{Reflejar y componer}\label{sec:reflejar}

{\small\vspace{1mm}\linkUnidad{03_reflexiones_similaridades}{Simulación interactiva: Unidad 5 — Reflexiones y similaridades}\par}


Una reflexión invierte una figura respecto de un centro o de una recta. Aplicada dos veces
respecto del mismo elemento devuelve la figura original. Esta propiedad se llama involución.

La riqueza geométrica aparece al combinar reflexiones distintas:
\begin{itemize}
  \item dos reflexiones puntuales producen una traslación;
  \item dos reflexiones en rectas paralelas producen una traslación perpendicular a ellas;
  \item dos reflexiones en rectas secantes producen una rotación alrededor de su intersección;
  \item tres reflexiones puntuales vuelven a producir una reflexión puntual y construyen un
  paralelogramo.
\end{itemize}

\paragraph{Primera traducción algebraica.}
\[
  \rho_b(x)=2b-x.
\]
La diferencia \(\rho_b(x)-b=-(x-b)\) indica que la imagen conserva la distancia a \(b\),
pero queda en la dirección opuesta. Al componer dos reflexiones, los términos que dependen de
\(x\) recuperan su signo y queda un desplazamiento constante.

\begin{idea}[Qué debe observarse]
Cada reflexión individual cambia la orientación. Dos reflexiones vuelven a preservarla. La
posición relativa de los centros o de los ejes decide si el resultado se interpreta como
traslación o como rotación.
\end{idea}

\begin{idea}[En palabras del profesor]
``Vimos que con reflexiones podíamos construir rotaciones (con un argumento coloquial
sobre que un par de reflexiones sucesivas no altera las distancias, pero tiene un punto
fijo, así que tiene que ser una rotación)'' (resumen de clase 3, cursada 2025).
\end{idea}

\subsection{Escalar: homotecias y similaridades}\label{sec:escalar}

Una homotecia cambia el tamaño respecto de un centro. Cada punto permanece sobre la recta que
lo une con ese centro. Si la razón es positiva conserva el semirayo; si es negativa pasa al
semirayo opuesto.

Una similaridad combina escala uniforme, giro y, eventualmente, reflexión o traslación. La
figura puede cambiar de tamaño y posición, pero conserva sus ángulos y sus proporciones. Dos
triángulos relacionados por una similaridad tienen la misma forma.

\paragraph{Primera traducción algebraica.}
\[
  x'=o+k(x-o).
\]
El término \(x-o\) mira el punto desde el centro. Multiplicarlo por \(k\) cambia su distancia
y, si \(k<0\), también invierte el semirayo.

\begin{idea}[En palabras del profesor]
``Así como la rotación se expresa en complejos también se expresa la escala, que en
matrices se representa con una matriz múltiplo de la identidad'' (resumen de clase 1, cursada 2025).
\end{idea}

\section{Construcciones, cuentas y ejemplos}

Este nivel es la guía de transformaciones. Reproducimos los enunciados de \texttt{Ejercicios.pdf}, material docente enviado por la cátedra, y resolvemos cada definición y cada teorema de dos maneras, con una construcción y con un ejemplo con vectores. La resolución es desarrollo nuestro.

{\small\vspace{1mm}\linkUnidad{03_reflexiones_similaridades}{Simulación interactiva: Unidad 5 — Reflexiones y similaridades}\par}

\subsection{Enunciados originales de la guía}

Los textos siguientes reproducen los enunciados de \texttt{Ejercicios.pdf}. Se normalizaron
únicamente los cortes de línea, la notación matemática y errores tipográficos evidentes; las
dos páginas originales se incluyen completas como anexo facsimilar al final del documento.

\begin{enunciado}[Definición 1]
La reflexión del punto \(a\) por un punto \(b\) es la operación que lleva el punto \(a\) al
punto que está sobre la recta \(\overline{ab}\), a la misma distancia de \(b\) que \(a\),
pero al otro lado de \(b\). Se llama \(\rho_b(a)\).
\end{enunciado}

\begin{enunciado}[Teorema 1]
\(\rho_a\rho_b\), la composición de dos reflexiones, es una traslación. Puede ser expresada
algebraicamente con la suma de un vector constante o geométricamente como el desplazamiento
según un segmento fijo.
\end{enunciado}

\begin{enunciado}[Teorema 2]
Si \(\rho_d=\rho_a\rho_b\rho_c\), entonces el polígono \(abcda\) es un paralelogramo.
\end{enunciado}

\begin{enunciado}[Observación]
Con esta operación puede definirse la suma y resta de vectores. ¿De qué modo? ¿Cómo se
relaciona esto con la propiedad de que sea inversiva la transformación, es decir, que al
aplicarla dos veces se anula?
\end{enunciado}

\begin{enunciado}[Definición 2]
La reflexión del punto \(a\) por la recta \(B\), \(\varphi_B(a)\), es la operación que lleva
el punto \(a\) al punto que está sobre la perpendicular a \(B\) que pasa por \(a\), a la misma
distancia que \(a\) de \(B\), pero al otro lado de la recta \(B\).
\end{enunciado}

\begin{enunciado}[Teorema 3]
Si \(\Theta_{AB}=\varphi_A\varphi_B\) y \(A\) es paralela a \(B\), entonces
\(\Theta_{AB}\) es una traslación. Puede ser expresada algebraicamente con la suma de un
vector constante o geométricamente como el desplazamiento según un segmento fijo.
\end{enunciado}

\begin{enunciado}[Teorema 4]
Si \(\Theta_{AB}=\varphi_A\varphi_B\) y \(A\) no es paralela a \(B\), entonces
\(\Theta_{AB}\) es una rotación.
\end{enunciado}

\begin{enunciado}[Teorema 5]
Si \(\Theta_{ABCD}=\Theta_{AB}\Theta_{CD}\), con \(\Theta_{AB}\) una traslación y
\(\Theta_{CD}\) una rotación, entonces \(\Theta_{ABCD}\) es una rotación. Además,
\(\Theta_{CDAB}=\Theta_{CD}\Theta_{AB}\) es también una rotación, pero es distinta.
\end{enunciado}

\begin{enunciado}[Definición 3]
Una homotecia con centro en el punto \(o\) y razón \(k\), nombrada \(\eta_o^k\), lleva un
punto \(a\) a otro sobre la recta \(\overline{oa}\), de forma que
\[
  \frac{d\bigl(o,\eta_o^k(a)\bigr)}{d(o,a)}=k,
\]
donde \(d(a,b)\) es la distancia entre el punto \(a\) y el punto \(b\). Los puntos \(a\) y
\(\eta_o^k(a)\) están al mismo lado del punto \(o\) en la recta \(\overline{oa}\).
\end{enunciado}

\begin{atencion}[Precisión sobre el texto original]
La razón entre distancias es siempre no negativa. Por eso, la redacción original corresponde
al caso \(k>0\). Para una razón firmada debe escribirse
\[
\eta_o^k(a)=o+k(a-o),
\qquad
\frac{d\bigl(o,\eta_o^k(a)\bigr)}{d(o,a)}=|k|.
\]
Punto e imagen están del mismo lado de \(o\) si \(k>0\), y en lados opuestos si \(k<0\).
\end{atencion}

\begin{enunciado}[Teorema 6]
Cualquier transformación de similaridad del plano puede expresarse como la composición
\(\varphi_A\varphi_B\eta_o^k\), donde el punto \(o\) es la intersección de las rectas
\(A\) y \(B\), y es el punto fijo de la transformación. Dado un triángulo y su imagen bajo
esta transformación puede hallarse el punto \(o\) con la construcción de Euler indicada en
la guía original.
\end{enunciado}

\begin{atencion}[Alcance del Teorema 6 original]
Tal como está escrito, el enunciado es demasiado general. La composición de dos reflexiones
y una homotecia describe una similaridad directa con punto fijo. Una similaridad inversa
requiere una cantidad impar de reflexiones, y una traslación no posee punto fijo finito.
\end{atencion}

\begin{enunciado}[Ejercicios]
Para cada una de las definiciones y teoremas expuestos, verificar la conclusión con un ejemplo
gráfico y construyendo un ejemplo algebraico con vectores. Pista para las reflexiones: para
hallar la proyección perpendicular de un punto sobre una recta puede usarse el producto
interno o construir la ecuación vectorial de la perpendicular usando el vector normal que
surge de la ecuación canónica de la recta.
\end{enunciado}

\subsection{Resolución completa de la guía}

La consigna de la guía no contiene una lista de problemas independientes. Solicita verificar
\textbf{cada definición y cada teorema} de dos maneras: mediante una construcción geométrica
y mediante un ejemplo algebraico con vectores. Las secciones siguientes constituyen esa
resolución punto por punto.

\begin{center}
\begin{tabularx}{\textwidth}{>{\bfseries}p{34mm} X X}
\toprule
Ítem de la guía & Verificación geométrica & Verificación algebraica\\
\midrule
Definición 1 & punto medio y segmentos iguales & \(\rho_b(x)=2b-x\)\\
Teorema 1 & desplazamiento constante & composición de dos reflexiones puntuales\\
Teorema 2 & diagonales que se bisecan & \(d=a-b+c\)\\
Definición 2 & perpendicular y distancias iguales & proyección mediante el vector normal\\
Teorema 3 & avance perpendicular de longitud \(2h\) & composición con normales paralelos\\
Teorema 4 & ángulo final igual a \(2\theta\) & producto de matrices de reflexión\\
Teorema 5 & mismo giro, distinto centro & resolución del punto fijo\\
Definición 3 & rayos desde el centro & \(\eta_o^k(x)=o+k(x-o)\)\\
Teorema 6 & escala y giro con centro común & \(S(x)=o+\lambda R(x-o)\)\\
\bottomrule
\end{tabularx}
\end{center}

\begin{idea}[Cómo reconocer que un ítem está resuelto]
Cada apartado presenta el dibujo o construcción, la propiedad geométrica que debe observarse,
la demostración vectorial y un ejemplo numérico. Los ejercicios del final son práctica
adicional y no sustituyen estas soluciones.
\end{idea}

\subsection{Guía resuelta: reflexión respecto de un punto}

\subsubsection{Definición 1}

La reflexión de \(x\) respecto del punto \(b\) deja a \(b\) como punto medio del segmento
que une \(x\) con su imagen. Por lo tanto,
\[
  \frac{x+\rhoP{b}(x)}{2}=b
  \qquad\Longrightarrow\qquad
  \boxed{\rhoP{b}(x)=2b-x}.
\]

\begin{center}
\begin{tikzpicture}[scale=1.0]
  \coordinate (x) at (1,1.5);
  \coordinate (b) at (5,2.4);
  \coordinate (xp) at (9,3.3);
  \draw[auxline] (.4,1.36)--(9.6,3.44);
  \draw[decorate,decoration={brace,amplitude=5pt,mirror},teal]
    (x)--node[below=5pt,font=\small] {$d(x,b)$}(b);
  \draw[decorate,decoration={brace,amplitude=5pt,mirror},teal]
    (b)--node[below=5pt,font=\small] {$d(b,\rho_b(x))$}(xp);
  \fill[teal] (x) circle (2.4pt) node[above left] {$x$};
  \fill[ink] (b) circle (2.7pt) node[above] {$b$};
  \fill[rust] (xp) circle (2.4pt) node[above right] {$\rho_b(x)$};
\end{tikzpicture}
\end{center}

\begin{resultado}[Propiedades inmediatas]
\[
  \rhoP{b}\bigl(\rhoP{b}(x)\bigr)=x,
  \qquad
  \rhoP{b}(b)=b,
  \qquad
  \rhoP{b}(x)-b=-(x-b).
\]
Es una transformación involutiva: aplicarla dos veces devuelve el punto original.
\end{resultado}

\paragraph{Ejemplo.}
Si \(b=(2,1)\) y \(x=(5,3)\), entonces
\[
  \rhoP{b}(x)=2(2,1)-(5,3)=(-1,-1).
\]
El punto medio entre \((5,3)\) y \((-1,-1)\) es efectivamente \((2,1)\).

\subsection{Guía resuelta: dos reflexiones puntuales}

\subsubsection{Teorema 1: la composición es una traslación}

Sean \(a,b\in\R^2\). Aplicando primero \(\rhoP{b}\) y luego \(\rhoP{a}\):
\[
\begin{aligned}
  \rhoP{a}\rhoP{b}(x)
  &=\rhoP{a}(2b-x)\\
  &=2a-(2b-x)\\
  &=x+2(a-b).
\end{aligned}
\]

\begin{resultado}[Conclusión]
\[
  \boxed{\rhoP{a}\rhoP{b}=T_{\,2(a-b)}}
\]
La composición desplaza todos los puntos por el mismo vector \(2(a-b)\). No posee puntos
fijos salvo el caso degenerado \(a=b\), en el cual es la identidad.
\end{resultado}

\begin{center}
\begin{tikzpicture}[scale=.92]
  \coordinate (x) at (.7,1);
  \coordinate (b) at (3,1.6);
  \coordinate (y) at (5.3,2.2);
  \coordinate (a) at (6.8,2.6);
  \coordinate (z) at (8.3,3);
  \draw[auxline] (.2,.87)--(8.8,3.13);
  \draw[mainline,-{Latex[length=2.8mm]}] (x)--(z);
  \fill[teal] (x) circle (2.4pt) node[labelback,below left] {$x$};
  \fill[ink] (b) circle (2.4pt) node[labelback,above] {$b$};
  \fill[rust] (y) circle (2.4pt) node[labelback,below] {$\rho_b(x)$};
  \fill[ink] (a) circle (2.4pt) node[labelback,above] {$a$};
  \fill[teal] (z) circle (2.4pt) node[labelback,above right] {$\rho_a\rho_b(x)$};
  \node[labelback,font=\small,text=tealdark,above=16pt] at ($(x)!.5!(z)$)
    {mismo desplazamiento para cualquier \(x\)};
\end{tikzpicture}
\end{center}

\paragraph{Ejemplo.}
Con \(a=(3,2)\) y \(b=(1,1)\), la traslación resultante tiene vector
\[
  2(a-b)=2(2,1)=(4,2).
\]
Para \(x=(2,-1)\): \(\rhoP{b}(x)=(0,3)\) y
\(\rhoP{a}(0,3)=(6,1)=x+(4,2)\).

\subsection{Guía resuelta: tres reflexiones puntuales}

\subsubsection{Teorema 2: aparece un paralelogramo}

La composición de tres reflexiones puntuales vuelve a ser una reflexión puntual:
\[
\begin{aligned}
\rhoP{a}\rhoP{b}\rhoP{c}(x)
  &=2a-\bigl(2b-(2c-x)\bigr)\\
  &=2(a-b+c)-x.
\end{aligned}
\]
Si \(\rhoP{d}=\rhoP{a}\rhoP{b}\rhoP{c}\), entonces
\[
  \boxed{d=a-b+c}.
\]
Al reordenar:
\[
  a+c=b+d.
\]
Esta es precisamente la condición vectorial para que las diagonales del cuadrilátero
\(a b c d\) tengan el mismo punto medio.

\begin{center}
\begin{tikzpicture}[scale=1.0]
  \coordinate (a) at (1,1);
  \coordinate (b) at (5,1.2);
  \coordinate (c) at (7.2,4);
  \coordinate (d) at (3.2,3.8);
  \draw[mainline] (a)--(b)--(c)--(d)--cycle;
  \draw[auxline] (a)--(c);
  \draw[auxline] (b)--(d);
  \fill[teal] (a) circle (2.4pt) node[below left] {$a$};
  \fill[teal] (b) circle (2.4pt) node[below right] {$b$};
  \fill[rust] (c) circle (2.4pt) node[above right] {$c$};
  \fill[rust] (d) circle (2.4pt) node[above left] {$d=a-b+c$};
  \fill[ink] (4.1,2.5) circle (2pt);
  \node[labelback,font=\small] at (4.15,2.02) {punto medio común};
  \draw[-{Latex[length=2.2mm]},rust,line width=.8pt] (a)--(b);
  \draw[-{Latex[length=2.2mm]},rust,line width=.8pt] (d)--(c);
\end{tikzpicture}
\end{center}

\paragraph{Ejemplo.}
Si \(a=(1,1)\), \(b=(4,1)\) y \(c=(6,3)\), entonces
\[
  d=(1,1)-(4,1)+(6,3)=(3,3).
\]
Se verifica \(a+c=(7,4)=b+d\); por lo tanto, las diagonales se bisecan.

\begin{idea}[Suma y resta de vectores]
La fórmula \(d=a-b+c\) permite construir geométricamente la operación
\(a+c-b\). Elegido un origen \(b\), el punto \(d\) representa la suma de los
vectores \(\overrightarrow{ba}+\overrightarrow{bc}\).

La suma \(u+v\) se obtiene llevando ambos vectores a un origen común y completando el
paralelogramo. La resta \(u-v\) se interpreta como \(u+(-v)\), donde \(-v\) surge de reflejar
el extremo de \(v\) respecto de su origen. Esto se relaciona con la involución porque
\[
  \rho_b^2=\operatorname{Id};
\]
reflejar dos veces no ``anula'' el punto, sino que devuelve exactamente la posición inicial.
\end{idea}

\subsection{Guía resuelta: reflexión respecto de una recta}

\subsubsection{Definición 2 y fórmula con proyección}

Sea la recta
\[
  B=\{x\in\R^2:n\cdot x=c\},
\]
donde \(n\neq0\) es un vector normal. La proyección ortogonal de \(x\) sobre \(B\) es
\[
  p=x-\frac{n\cdot x-c}{\lVert n\rVert^2}\,n.
\]
Como \(p\) debe ser el punto medio entre \(x\) y su reflejo:
\[
  \boxed{\refl{B}(x)=2p-x
  =x-2\frac{n\cdot x-c}{\lVert n\rVert^2}\,n}.
\]

\begin{center}
\begin{tikzpicture}[scale=1.0]
  \coordinate (x) at (2,4.4);
  \coordinate (p) at (4.8,2.4);
  \coordinate (xp) at (7.6,.4);
  \draw[rust,line width=1.1pt] (.4,.45)--(8.8,3.95)
    node[right,font=\small] {$B$};
  \draw[auxline] (x)--(xp);
  \draw[mainline] ($(p)+(-.32,.15)$)--($(p)+(-.17,.47)$)--($(p)+(.15,.32)$);
  \fill[teal] (x) circle (2.5pt) node[above left] {$x$};
  \fill[ink] (p) circle (2.4pt) node[above right] {$p$};
  \fill[rust] (xp) circle (2.5pt) node[below right] {$\varphi_B(x)$};
  \draw[decorate,decoration={brace,amplitude=5pt},teal]
    (x)--node[left=5pt,font=\small] {$d$}(p);
  \draw[decorate,decoration={brace,amplitude=5pt},teal]
    (p)--node[left=5pt,font=\small] {$d$}(xp);
\end{tikzpicture}
\end{center}

\paragraph{Ejemplo.}
Para \(B:\ x+2y=4\), \(n=(1,2)\) y \(x_0=(4,3)\):
\[
  n\cdot x_0-c=6,\qquad \lVert n\rVert^2=5,
\]
\[
  \refl{B}(x_0)
  =(4,3)-\frac{12}{5}(1,2)
  =\left(\frac85,-\frac95\right).
\]
La proyección es
\[
  p=(4,3)-\frac65(1,2)
  =\left(\frac{14}{5},\frac35\right).
\]
Además,
\[
  x_0-p=p-\refl{B}(x_0)=\frac65(1,2),
\]
lo que verifica simultáneamente que ambos segmentos tienen la misma longitud y que son
perpendiculares a la recta, pues son paralelos al normal \(n=(1,2)\).

\subsection{Guía resuelta: reflexiones en rectas paralelas}

\subsubsection{Teorema 3: la composición es una traslación}

Tomemos rectas paralelas con un normal unitario \(n\):
\[
  A:\ n\cdot x=\alpha,
  \qquad
  B:\ n\cdot x=\beta.
\]
Sus reflexiones son
\[
  \refl{A}(x)=x-2(n\cdot x-\alpha)n,
  \qquad
  \refl{B}(x)=x-2(n\cdot x-\beta)n.
\]
Al componer:
\[
  \boxed{\refl{A}\refl{B}(x)=x+2(\alpha-\beta)n}.
\]

\begin{center}
\begin{tikzpicture}[scale=.94]
  \draw[rust,line width=1.1pt] (1,.5)--(1,5) node[above] {$B$};
  \draw[rust,line width=1.1pt] (4,.5)--(4,5) node[above] {$A$};
  \coordinate (x) at (0,1.2);
  \coordinate (y) at (2,1.2);
  \coordinate (z) at (6,1.2);
  \draw[auxline] (x)--(z);
  \draw[mainline,-{Latex[length=2.8mm]}] (x)--(z);
  \fill[teal] (x) circle (2.5pt) node[below] {$x$};
  \fill[rust] (y) circle (2.5pt) node[below] {$\varphi_B(x)$};
  \fill[teal] (z) circle (2.5pt) node[below] {$\varphi_A\varphi_B(x)$};
  \draw[<->,muted] (1,4.1)--node[above,font=\small] {$h$}(4,4.1);
  \node[font=\small,text=tealdark] at (3,2.3) {desplazamiento \(2h\)};
\end{tikzpicture}
\end{center}

\paragraph{Lectura geométrica.}
La primera reflexión cruza \(B\); la segunda cruza \(A\). El desplazamiento total es
perpendicular a ambas rectas y su módulo es dos veces la distancia entre ellas. Invertir el
orden invierte el vector de traslación.

\paragraph{Ejemplo gráfico y algebraico común.}
Si \(B:x=1\), \(A:x=4\) y \(p=(0,2)\), entonces
\[
  p=(0,2)
  \xmapsto{\ \refl{B}\ }(2,2)
  \xmapsto{\ \refl{A}\ }(6,2).
\]
Por tanto,
\[
  \refl{A}\refl{B}(x,y)=(x+6,y).
\]
La distancia entre los ejes es \(3\), la traslación tiene módulo \(6\) y el punto intermedio
\((2,2)\) muestra las dos reflexiones separadamente.

\subsection{Guía resuelta: reflexiones en rectas secantes}

\subsubsection{Teorema 4: la composición es una rotación}

Sean \(A\) y \(B\) dos rectas que se cortan en \(o\), formando un ángulo orientado
\(\theta\) desde \(B\) hasta \(A\). Como ambas reflexiones fijan \(o\), su composición
también lo fija. Además, cada reflexión invierte orientación; dos reflexiones la preservan.
El efecto angular total es
\[
  \boxed{\refl{A}\refl{B}=R_{o,\,2\theta}}.
\]

\begin{center}
\begin{tikzpicture}[scale=1.0]
  \coordinate (o) at (4,2.2);
  \coordinate (x) at (7.4,1.35);
  \coordinate (y) at (6.5,4.1);
  \draw[rust,line width=1.05pt] (.4,2.2)--(8,2.2) node[right] {$B$};
  \draw[rust,line width=1.05pt] ($(o)+(-3,-1.73)$)--($(o)+(3,1.73)$)
    node[right] {$A$};
  \draw[auxline] (o)--(x);
  \draw[auxline] (o)--(y);
  \draw[mainline,-{Latex[length=2.6mm]}] (x) arc[start angle=-14,end angle=39,
    radius=3.5];
  \fill[ink] (o) circle (2.6pt) node[below left] {$o$};
  \fill[teal] (x) circle (2.5pt) node[below right] {$x$};
  \fill[teal] (y) circle (2.5pt)
    node[labelback,above right=2pt] {$\varphi_A\varphi_B(x)$};
  \draw[teal,line width=.7pt] ($(o)+(.9,0)$)
    arc[start angle=0,end angle=30,radius=.9];
  \node[labelback,font=\small,text=teal] at ($(o)+(1.15,.32)$) {$\theta$};
  \node[labelback,font=\small,text=tealdark] at (6.05,2.86) {giro total \(2\theta\)};
\end{tikzpicture}
\end{center}

\paragraph{Verificación matricial.}
Si ambas rectas pasan por el origen y \(A\) forma ángulo \(\alpha\), la matriz de reflexión es
\[
  H(\alpha)=
  \begin{pmatrix}
    \cos 2\alpha & \sin 2\alpha\\
    \sin 2\alpha &-\cos 2\alpha
  \end{pmatrix}.
\]
Por multiplicación directa,
\[
  H(\alpha)H(\beta)=
  \begin{pmatrix}
    \cos 2(\alpha-\beta)&-\sin 2(\alpha-\beta)\\
    \sin 2(\alpha-\beta)& \cos 2(\alpha-\beta)
  \end{pmatrix}
  =R\bigl(2(\alpha-\beta)\bigr).
\]

\paragraph{Ejemplo numérico completo.}
Sea \(B\) el eje \(x\), sea \(A\) una recta por el origen que forma \(30^\circ\) con \(B\),
y sea \(p=(1,1)\). Primero:
\[
  p_1=\varphi_B(p)=(1,-1).
\]
La reflexión respecto de \(A\) tiene matriz
\[
H(30^\circ)=
\begin{pmatrix}
\frac12&\frac{\sqrt3}{2}\\
\frac{\sqrt3}{2}&-\frac12
\end{pmatrix}.
\]
Por lo tanto,
\[
p_2=\varphi_A(p_1)
=
\begin{pmatrix}
\frac12-\frac{\sqrt3}{2}\\[1mm]
\frac{\sqrt3}{2}+\frac12
\end{pmatrix}.
\]
Este es exactamente
\[
R(60^\circ)(1,1)^T.
\]
El punto \(p\) se refleja primero debajo de \(B\) y luego cruza \(A\); el ángulo neto observado
desde el origen es el doble del ángulo entre los ejes.

\begin{idea}[Puente con la primera clase]
Esta matriz es exactamente la matriz de rotación trabajada en clase. La guía aporta ahora
una construcción geométrica: una rotación plana puede generarse con dos reflexiones en
rectas que se cortan en su centro.
\end{idea}

\subsection{Guía resuelta: traslación y rotación}

\subsubsection{Teorema 5: el orden cambia el centro}

Sea \(T_v(x)=x+v\) y sea una rotación de parte lineal \(R\) y centro \(c\):
\[
  Q_c(x)=c+R(x-c)=Rx+(I-R)c.
\]

\paragraph*{Primero rotar, después trasladar}
\[
  T_vQ_c(x)=Rx+(I-R)c+v.
\]
Si el ángulo no es \(0\) módulo \(2\pi\), la matriz \(I-R\) es invertible y existe un único
centro \(c_1\):
\[
  \boxed{c_1=c+(I-R)^{-1}v}.
\]

\paragraph*{Primero trasladar, después rotar}
\[
  Q_cT_v(x)=Rx+(I-R)c+Rv,
\]
por lo que su centro es
\[
  \boxed{c_2=c+(I-R)^{-1}Rv}.
\]
En general \(c_1\neq c_2\): ambas composiciones tienen el mismo ángulo, pero centros distintos.

\begin{center}
\begin{tikzpicture}[scale=.92]
  \coordinate (c) at (4,2.4);
  \coordinate (c1) at (2,4.4);
  \coordinate (c2) at (6.4,4.1);
  \coordinate (x) at (5.8,1);
  \draw[auxline] (c)--(x);
  \draw[-{Latex[length=2.5mm]},teal,line width=1pt] (x) arc[start angle=-35,end angle=35,radius=2.2];
  \draw[-{Latex[length=2.5mm]},rust,line width=1pt]
    (1.3,.8)--node[labelback,below,font=\small] {$v$}(3.1,1.45);
  \fill[ink] (c) circle (2.6pt) node[labelback,below] {$c$};
  \fill[teal] (c1) circle (2.6pt) node[labelback,above left] {$c_1$};
  \fill[rust] (c2) circle (2.6pt) node[labelback,above right] {$c_2$};
  \draw[teal,dashed] (c1) circle (1.25);
  \draw[rust,dashed] (c2) circle (1.25);
  \node[labelback,font=\small,text=teal] at (1.85,3.15) {\(T_vQ_c\)};
  \node[labelback,font=\small,text=rust] at (7.05,3.0) {\(Q_cT_v\)};
\end{tikzpicture}
\end{center}

\begin{atencion}[Caso excepcional]
Si \(R=I\), no hay una rotación no trivial: las composiciones son simplemente la traslación
\(T_v\). Por eso el enunciado debe entenderse para una rotación de ángulo no nulo.
\end{atencion}

\paragraph{Ejemplo de orden.}
Sea \(R\) el giro de \(90^\circ\) alrededor del origen y \(v=(2,0)\). Entonces
\[
  T_vR(x)=Rx+(2,0),
  \qquad
  RT_v(x)=Rx+(0,2).
\]
El centro del primero es \((1,1)\); el del segundo es \((-1,1)\). El ángulo es el mismo,
pero las transformaciones son distintas.

Siguiendo el mismo punto \(p=(1,0)\):
\[
p\xmapsto{\ R\ }(0,1)
\xmapsto{\ T_v\ }(2,1),
\]
mientras que
\[
p\xmapsto{\ T_v\ }(3,0)
\xmapsto{\ R\ }(0,3).
\]
La separación entre \((2,1)\) y \((0,3)\) permite ver directamente que el orden cambió la
transformación.

\subsection{Guía resuelta: homotecia}

\subsubsection{Definición 3}

La homotecia de centro \(o\) y razón \(k\) se escribe
\[
  \boxed{\homot{o}{k}(x)=o+k(x-o)=kx+(1-k)o}.
\]

\begin{center}
\begin{tikzpicture}[scale=.95]
  \coordinate (o) at (1,1);
  \coordinate (a) at (4,2.2);
  \coordinate (ap) at (7,3.4);
  \draw[auxline] (.5,.8)--(8,3.8);
  \draw[mainline,-{Latex[length=2.6mm]}] (a)--(ap);
  \fill[ink] (o) circle (2.6pt) node[below left] {$o$};
  \fill[teal] (a) circle (2.6pt) node[above left] {$a$};
  \fill[rust] (ap) circle (2.6pt) node[above right] {$\eta_o^k(a)$};
  \draw[decorate,decoration={brace,amplitude=5pt,mirror},teal]
    (o)--node[below=5pt,font=\small] {$r$}(a);
  \draw[decorate,decoration={brace,amplitude=5pt,mirror},rust]
    (o)--node[below=5pt,font=\small] {$kr$}(ap);
\end{tikzpicture}
\end{center}

\begin{center}
\begin{tabularx}{.92\textwidth}{>{\bfseries}c X}
\toprule
Valor de \(k\) & Efecto geométrico\\
\midrule
\(k>1\) & ampliación; imagen y punto están del mismo lado de \(o\)\\
\(0<k<1\) & reducción; imagen entre \(o\) y el punto\\
\(k=1\) & identidad\\
\(k=0\) & todos los puntos colapsan en \(o\); no es una similaridad invertible\\
\(k<0\) & escala por \(|k|\) y cambia al semirayo opuesto\\
\bottomrule
\end{tabularx}
\end{center}

\paragraph{Ejemplo.}
Con \(o=(1,1)\), \(k=2\) y \(a=(3,2)\):
\[
  \homot{o}{2}(a)=(1,1)+2(2,1)=(5,3).
\]
La distancia al centro se duplica y la dirección desde \(o\) se conserva.

\subsection{Guía resuelta: similaridades}

\subsubsection{Teorema 6}

Una similaridad conserva ángulos y multiplica todas las distancias por un mismo factor
\(\lambda>0\). Una similaridad directa con punto fijo \(o\) tiene la forma
\[
  S(x)=o+\lambda R(x-o),
\]
donde \(R\) es una rotación. Como toda rotación plana puede expresarse como composición de
dos reflexiones en rectas que se cortan en \(o\),
\[
  \boxed{S=\refl{A}\refl{B}\homot{o}{\lambda}}.
\]
La homotecia conmuta con las reflexiones cuyos ejes pasan por \(o\), porque todas actúan
sobre el vector \(x-o\).

\begin{center}
\begin{tikzpicture}[scale=.88]
  \coordinate (o) at (1.3,1.1);
  \coordinate (a) at (3.1,1.4);
  \coordinate (b) at (2.2,3);
  \coordinate (c) at (4,2.7);
  \coordinate (ap) at (6.2,2.1);
  \coordinate (bp) at (5.7,5.0);
  \coordinate (cp) at (8.6,4.3);
  \draw[teal,line width=1.1pt] (a)--(b)--(c)--cycle;
  \draw[rust,line width=1.1pt] (ap)--(bp)--(cp)--cycle;
  \draw[auxline] (o)--(a);
  \draw[auxline] (o)--(ap);
  \fill[ink] (o) circle (2.7pt) node[below left] {$o$};
  \foreach \p/\lab in {a/a,b/b,c/c}
    \fill[teal] (\p) circle (2.1pt)
      node[labelback,above left=1.5pt,font=\small] {$\lab$};
  \foreach \p/\lab in {ap/S(a),bp/S(b),cp/S(c)}
    \fill[rust] (\p) circle (2.1pt)
      node[labelback,above right=1.5pt,font=\small] {$\lab$};
  \draw[-{Latex[length=2.6mm]},tealdark,line width=.9pt] (3.9,.7)
    .. controls (4.8,.2) and (5.6,.8) .. (6.1,1.5);
  \node[font=\small,text=tealdark] at (5,.35) {escala \(+\) giro};
\end{tikzpicture}
\end{center}

\begin{atencion}[Alcance del enunciado]
La composición de dos reflexiones preserva la orientación. Por eso la fórmula anterior
describe las similaridades directas. Una similaridad inversa requiere una cantidad impar de
reflexiones, por ejemplo una reflexión y una homotecia. Si la similaridad no tiene punto fijo
finito, aparece el caso especial de una traslación o una transformación con componente de
traslación.
\end{atencion}

\section{En profundidad}

Cerramos el tema con dos complementos de la guía: un ejemplo numérico completo de la descomposición de una similaridad y el método para recuperar el centro a partir de un triángulo y su imagen. Ambos son desarrollo nuestro.

\subsection{Ejemplo numérico de la descomposición}

Tomemos \(o=(0,0)\), una escala \(\lambda=2\) y una rotación de \(90^\circ\). Para el
triángulo
\[
  a=(1,0),\qquad b=(0,1),\qquad c=(1,1),
\]
la similaridad
\[
  S(x)=2R(90^\circ)x
\]
produce
\[
  a'=(0,2),\qquad b'=(-2,0),\qquad c'=(-2,2).
\]
Elija como primer eje de reflexión \(B\) el eje \(x\), y como segundo eje \(A\) la recta
\(y=x\), que forma \(45^\circ\) con \(B\). Entonces
\[
  \varphi_A\varphi_B=R(90^\circ),
\]
porque el giro resultante es el doble del ángulo entre los ejes. Finalmente, la homotecia
\(\eta_o^2\) duplica todas las distancias al origen:
\[
  \boxed{S=\varphi_A\varphi_B\eta_o^2}.
\]
Geométricamente se ven tres etapas: reflexión en \(B\), reflexión en \(A\) y expansión radial
desde \(o\).

\subsection{Cómo hallar el centro a partir de un triángulo y su imagen}

Supongamos que una similaridad directa lleva los vértices
\[
  a\mapsto a',\qquad b\mapsto b',\qquad c\mapsto c'.
\]
Primero se determina el factor de escala y el ángulo:
\[
  \lambda=\frac{\lVert b'-a'\rVert}{\lVert b-a\rVert},
  \qquad
  \theta=\angle\bigl(b-a,b'-a'\bigr).
\]
Con estos datos, el centro \(o\) satisface
\[
  a'=o+\lambda R(\theta)(a-o).
\]
Reordenando:
\[
  \bigl(I-\lambda R(\theta)\bigr)o
  =a'-\lambda R(\theta)a,
\]
y, cuando la matriz es invertible,
\[
  \boxed{
  o=
  \bigl(I-\lambda R(\theta)\bigr)^{-1}
  \left(a'-\lambda R(\theta)a\right)
  }.
\]
El mismo resultado debe obtenerse usando \(b\mapsto b'\) o \(c\mapsto c'\).

\paragraph{Construcción geométrica.}
Como la similaridad multiplica todas las distancias al centro por \(\lambda\),
\[
  \lVert o-a'\rVert=\lambda\lVert o-a\rVert,
  \qquad
  \lVert o-b'\rVert=\lambda\lVert o-b\rVert.
\]
El primer lugar geométrico es el círculo de Apolonio asociado al par \((a,a')\), y el segundo
es el correspondiente a \((b,b')\). Sus intersecciones proporcionan los centros candidatos.
Se elige el que verifica el ángulo orientado \(\theta\). Si \(\lambda=1\), ambos lugares se
reducen a las mediatrices de \(aa'\) y \(bb'\), como ocurre en una rotación pura.

\begin{resultado}[Verificación del Teorema 6]
El centro se obtiene como intersección de dos condiciones geométricas de razón de distancias.
Una vez hallado, todos los segmentos que parten de \(o\) se escalan por \(\lambda\) y giran
el mismo ángulo \(\theta\). Esto completa la parte constructiva y la parte algebraica pedidas
por la guía.
\end{resultado}
